% Folder 83, batch 03 (archivists' pages 41-60).
% Pass: Opus 5.5 (claude-opus-5-5), 2026-10-10 — first pass, unchecked against the pages by a human.
%
% Page 46 is blank and page 51 is a cover sheet in his hand (its words head the
% section before page 52); neither has a \page{}. Pages 41-45 carry his own
% numbering 4-8: the run starts in the previous batch. The hand is fast and the
% connective prose is often lost; the formulas are read with more confidence.
\documentclass[11pt,a4paper]{article}
\input{../preamble/grothendieck}

\folder{83}
\batch{3}
\pages{41}{60}
\foldertitle{[Icosaèdres, bi-icosaèdres et algèbres d'Azumaya de rang 4] : notes manuscrites (1982-1983, 1986, s.d.).}
\dating{1982-1986}
\watermark{Édition de démonstration}

\begin{document}

\page{41}
\note{p. 4 de l'auteur. L'argument vient du lot précédent (pp. 1--3 de l'auteur) : il y est question d'objets $V_i$ indexés par un ensemble $I$ à quatre éléments, d'isomorphismes $u$ entre eux et d'un sous-objet $V$ de $\prod V_i$ ; les relations (1) et la définition des $u$ ne sont pas sur ces pages.}

\[
\text{(2)}\qquad u_{(i_2,i_1)(i_3,i_4)} = u_{(i_1,i_2)(i_3,i_4)}^{-1}
\]
\marginal{12 relations, \uncertain{reste à fixer} 12 quantités}

\struck{$u_{(i_2,i_1)(i_3,\ldots)}$}\note{une première écriture de (3), biffée et en partie illisible}

(3) \struck{$u_{i_2 i_1}$} \add{$u_{(i_1,i_3)(i_2,i_4)}$}
\[
\text{(3)}\qquad u_{(i_3,i_2)(i_1,i_4)}\,u_{(i_2,i_1)(i_3,i_4)} = \mathrm{id}
\]
\note{le premier facteur de (3) est biffé ; le facteur $u_{(i_1,i_3)(i_2,i_4)}$ écrit au-dessus semble pris dans le même trait, et la lecture des indices des deux derniers facteurs est douteuse}

$V$ est connu en tant que sous-objet de $\prod V_i$, quand on connaît le système des $u$, comme étant formé des $(x_i)_{i\in I}$ tels que pour \struck{tout} $i\in I$ \uncertain{donné} les $x_j$ ($j\in I\smallsetminus\{i\}$) soient « de somme nulle » pour la structure \uncertain{tautologique} d'iso entre les $V_j$ ($j\in J$) définie par l'une ou l'autre (\uncertain{c'est la même} $=$) des deux permutations circulaires etc. Mais il faut des conditions pour que le $W$ ainsi défini (qui a priori serait tel que $W\to V_i\times V_j = V_J$ \struck{soit} soit injectif) soit « assez gros » pour que ce soit bijectif. On va écrire des conditions pour ceci.

\page{42}
\note{p. 5 de l'auteur}
On va définir, pour un système d'iso $u_{ijkl}$, un système de « rapports anharmoniques »
\[
\text{(4)}\qquad \rho_{ijkl} \in \mathrm{Aut}\,V_i \qquad (\{i,j,k,l\} = I).
\]
On a
\[
u_{ijkl} : V_i \xrightarrow{\ \sim\ } V_j, \qquad u_{ijlk} : V_i \xrightarrow{\ \sim\ } V_j,
\]
d'où
\[
\text{(4)}\qquad \rho_{ijkl} = u_{ijlk}^{-1}\,u_{ijkl}
\]
\note{le numéro (4) est répété, écrit sur un autre chiffre}
\marginal{24 quantités}

Donc on aura
\[
\text{(5)}\qquad \rho_{ijlk} = (\rho_{ijkl})^{-1}
\]
\marginal{\struck{12} relations, restent 12 quantités, correspondant aux \uncertain{arêtes orientées} de $I$}

On aura d'ailleurs
\[
\text{(6)}\qquad \rho_{jikl} = \text{\struck{$\ldots$}}
\begin{cases}
u_{ijkl}(\rho_{ijkl}^{-1})\ \text{\struck{$\ldots$}}\\
u_{ijlk}(\rho_{ijkl}^{-1})
\end{cases}
\qquad \rho_{jikl} = u_{jilk}^{-1}\,u_{jikl}
\]
\note{la seconde égalité, $\rho_{jikl}=u_{jilk}^{-1}u_{jikl}$, est écrite verticalement sous le membre de gauche}

N.B. On pose $\beta(\alpha) = \beta\alpha\beta^{-1}$.

[Ce qui donne
\[
\rho_{ji,lk} = u_{ijkl}(\rho_{ijkl}) = u_{ijlk}(\rho_{ijkl})
\]
\marginal{Invariance par double permutation}

\page{43}
\note{p. 6 de l'auteur}
Donc il reste \uncertain{à contrôler} \struck{\ill{}} 6 quantités, i.e. \struck{pour} il suffit, pour \uncertain{toute} \struck{partition (\ill{}) \ill{} de $I$ ; des $I'\sqcup I''$} partie à deux éléments $I'\subset I$, $I'=\{i'_1,i'_2\}$, $I''=\{i''_1,i''_2\}$, de déterminer l'une des quatre quantités
\[
\rho_{i'_1 i'_2\, i''_1 i''_2}
\]
\note{sous chacune des deux paires d'indices, une flèche double indique l'échange des deux éléments}
Nous allons \struck{\ill{}} donner une formule pour relier
\[
\rho_{ijkl} \quad\text{et}\quad \rho_{ikjl}
\]
(ce qui permettra de se passer d'un \uncertain{repère} : un \uncertain{repère} quelconque)
\[
\rho_{ijkl} + \rho_{ikjl} = 1
\]
i.e.
\[
\text{(7)}\qquad \boxed{\rho_{ikjl} = 1 - \rho_{ijkl}}
\]
\marginal{[Faire vérification}

\textbf{Prop.} Soit un système de $u_{ijkl} : V_i \to V_j$, satisfaisant les relations (2) (3) ; d'où (4) des $\rho_{ijkl}\in\mathrm{Aut}(V_i)$, satisfaisant (5)(6). Pour que \uncertain{ce} corresponde à une \uncertain{quadralité} \ill{} $(V,(V_i),(\varphi_i))$, \struck{\ill{}} $V\subset\prod V_i$, il

\page{44}
\note{p. 7 de l'auteur}
\uncertain{faut} et il suffit qu'on ait aussi les 24 relations (7). Quand \add{les} $\mathrm{End}(V_i)$ sont commutatifs (\uncertain{rappelons} qu'ils sont iso entre eux) \uncertain{donc} \struck{les \ill{}} \add{les seules} \struck{$\mathrm{Aut}\,V_i$} relations (2), (3) impliquent que les \add{\uncertain{syst.} \uncertain{transitifs}} d'iso entre les $\mathrm{End}(V_i)$, sont \struck{\ill{}} canoniques. Alors les relations 5, 6, 7 s'interprètent comme des relations dans \struck{\ill{}} $k$
\[
\text{(8)}\qquad
\begin{cases}
\rho_{jikl} = \rho_{ijlk} = \rho_{ijkl}^{-1}\\
\rho_{ikjl} = 1 - \rho_{ijkl}
\end{cases}
\qquad (\Rightarrow\ \rho_{jilk} = \rho_{ijkl})
\]
\note{le premier symbole de (8) est surchargé ; la lecture $\rho_{jikl}$ est douteuse}

En d'autres termes, faisons opérer $\mathfrak S_3$ sur l'ens. \add{$U(k)$} \struck{des} des $\lambda\in k$ tels que $\lambda$ et $1-\lambda$ soient inversibles, par
\[
\sigma_{12}(\lambda) = \sigma_{34}(\lambda) = \lambda^{-1}, \qquad \sigma_{23}(\lambda) = 1-\lambda
\]
\marginal{donc $\mathfrak S_4$ opère par l'intermédiaire du quotient $\mathfrak S_3$ de $\mathfrak S_4$}
l'application \struck{\ill{}} \struck{$\rho_{ijkl}$} $(i,j,k,l)\mapsto\rho_{ijkl}$,
\[
\underbrace{\mathrm{Rep}(I)}_{\text{torseur sous }\mathfrak S_4} \longrightarrow U(k)
\]
est compatible avec $\mathfrak S_4$.
\note{le nom de l'ensemble cible, ici et p. 45 (feuillet 45), est lu $U(k)$ sans certitude}

\page{45}
\note{p. 8 de l'auteur}
Autre façon de le dire : $\rho_{ijkl}$ ne dépend que du « coin » $(i,j)(k,l)$ muni d'un élément du produit contracté des \struck{\ill{}} codiagonales, i.e. d'une codiagonale, car
\[
\text{(9)}\qquad \rho_{(i,j)(k,l)} = \rho_{(k,l)(i,j)}
\]
\drawing{un carré dont les sommets sont joints par les deux diagonales en pointillé, marqué $Q_1$ et $Q_2$, avec une flèche vers $Q_1$ ; à côté : « On a le troisième coin »}
(N.B. Cette relation n'a de sens que parce que $\mathrm{End}(W)$ est supposé commutatif)

Si on se donne une \struck{coin} str. \uncertain{coin} sur une codiag. de $I$, on se donne \ill{} repère de l'ens. $\mathrm{Cor}(I)$ des structures coin, et la donnée de $\rho$ équivaut à une application de
\[
\underbrace{\mathrm{Rep}\,\mathrm{Cor}(I)}_{\text{torseur sous }\mathfrak S_3} \longrightarrow U(k)
\]
\marginal{vérifier \struck{\ill{}} \uncertain{la} \uncertain{compatibilité} ! il faut \ill{} \ill{} \ill{} $\mathfrak S_3$ …}

Il reste à décrire la catégorie des \uncertain{quadralités} dans $C$, \struck{\ill{}} de façon plus simple que par un système de 24 flèches $(u_{ijkl})$ données à deux \uncertain{isomorphismes} près.

\page{47}
\marginal{Avril 86}
\emph{Trialités.} Soit $C$ catégorie, $n$ entier naturel $\geq$ \uncertain{$2$}. On appelle \emph{$n$-alité} (trialité, quadralité, pentalité si $n=3,4,5$) un système
\marginal{(le cas $n=2$ trivial)}
\[
\{I, E, \{E_i\}_{i\in I}, (\varphi_i)_{i\in I}\}
\]
où $I$ \struck{ens.} de cardinal $n$,
\[
\varphi_i : E \to E_i \qquad i\in I
\]
flèches de $C$, satisfaisant :

$\forall J\subset I$ avec $\mathrm{card}\,J = 2$, l'application
\[
\text{(1)}\qquad E \xrightarrow{\ \varphi_J\ } E_J = \prod_{i\in J} E_i
\]
\struck{on suppose que ce produit existe}\note{biffé, sous le produit, qu'une accolade désigne}
est un \emph{iso} (dans $\widehat C$) $\forall J\in\mathfrak P_2(I)$, \uncertain{i.e.} $E$ muni des $\varphi_i$ ($i\in J$) est un produit des $E_i$ dans $C$.

[Les $n$-alités forment une catégorie $n\text{-al}(C)$, fibrée \add{sur} \struck{\ill{}} le groupoïde \uncertain{$\mathrm{Ens}_n$}.

Je m'intéresse surtout aux deux cas

1°) $C$ catégorie additive

2°) $C = \mathrm{Ens}$

\struck{Dans le cas 2} Je vais ici regarder le cas 2°). Si un des $E_i$ est vide, donc \add{$E$ l'est} (tous sauf au plus un le sont). On va exclure ce cas, \add{dit} « \struck{\ill{}} impropre ». Dans tous les autres cas les $E_i$ peuvent être considérés comme des quotients de $E$. Ainsi, la structure peut être considérée comme une partie \uncertain{par} $E$, \add{$I$} \struck{munie d'une}, i.e. une application $I\to\mathrm{Quot}(E)$. Cette application est d'ailleurs injective, sauf \struck{quand} \struck{\add{dans}}

\page{48}
le cas où tous les $E_i$ sont réduits à un élément \struck{\ill{}} (\struck{\ill{}} un des $E_i$ l'est $\Leftrightarrow$ $E$ l'est), qu'on dit « trivial ». Donc une $n$-alité non triviale peut être considérée comme un ens. $E$ muni d'une partie $I\subset\mathrm{Quot}(E)$, avec $\mathrm{card}\,I = n$, et satisfaisant l'axiome évident traduisant (1).

Autre version : Soit \struck{\ill{}}
\[
\text{(2)}\qquad S = \coprod_{i\in I} E_i
\]
(« sommets »)
et \struck{\ill{}} soit
\[
\text{(3)}\qquad R\subset S\times E
\]
défini par
\[
(s, x)\in R \iff s = \varphi_i(x) \qquad (s = (i\in I,\ s\in E_i),\ x\in E).
\]
\struck{Alors on a ceci}

Considérons
\[
E \longrightarrow \mathfrak P(S), \qquad x \longmapsto R(x)\subset S ;
\]
les \uncertain{parties} de \struck{$R(x)$} la forme $R(x)$ sont appelés $n$-uples de $S$. On dit que \struck{$E\to\mathfrak P(S)$ est \ill{}} $a\in\mathfrak P_2(S)$ est une \emph{arête} si elle est contenue dans un $n$-uple, i.e. $\{s,t\}$ est une arête ssi $s\neq t$ et $\exists x\in E$ tel que $\{s,t\}\subset R(x)$ i.e.

\page{49}
$(s,x), (t,x)\in R$. On a

(i) $\forall s,t\in S$, $s\neq t$, existe au plus un $x\in E$ tel que $(s,x),(t,x)\in R$ i.e. $s,t\in R(x)$ i.e. $\{s,t\}\subset R(x)$.

\struck{[N.B. cette condition est symétrique \ill{} $S$, $R$ \ill{} signifie \ill{} : si $s\neq t$, $x\neq y$ ($s,t\in S$, $x,y\in R$) alors il existe au plus un \ill{} $\{s,t\}$ \ill{}]}

\struck{(ii) La relation dans}

N.B. Cette \struck{relation} \add{condition} est symétrique en $S$, $R$. Elle signifie que si $a\in\mathfrak P_2(S)$, $u\in\mathfrak P_2(R)$, alors $\exists s\in a$, $x\in u$ tels que $(s,x)\notin R$.
\marginal{\uncertain{toute arête est contenue dans un seul $n$-uple}}

Si on suppose que les \struck{$R(x)$} $\mathrm{card}\,R(x)\geq 2$ \struck{\ill{}} (ce qui est le cas \add{auquel cas $\mathrm{card}\,R(x) = n$} \uncertain{pourvu} qu'on ait une \struck{\ill{}} $n$-alité, \struck{non impropre}), alors $R\to\mathfrak P(S)$ est injectif. De même, si on suppose que pour tt $s\in S$, $R(s)\subset R$ est de cardinal $\geq 2$ (ce qui est le cas si on part d'une $n$-alité propre non triviale) alors $S\to\mathfrak P(R)$ injectif)

(ii) La relation dans $S$ : « $\{s,t\}$ n'est pas une arête » (i.e. $s=t$ ou $\nexists\, x\in E$ tel que $\{s,t\}\subset R(x)$) est une relation d'équivalence dans $S$, i.e. pour une \struck{\ill{}} arête $\{s,t\}$, tout \uncertain{sommet} \struck{de} $S$ est lié soit à $s$, soit à $t$.

\struck{Cette} La condition (ii) implique que $S = \coprod E_i$ de telle façon que \struck{\ill{}} $a\in\mathfrak P_2(S)$, la est une

\page{50}
arête ssi $a$ n'est contenue dans aucun des \uncertain{$S_i$}. D'autre part, si $u$ est un $n$-uple, alors
\[
u \longrightarrow I \qquad (\text{i.e. } u\cap E_i \text{ de card} \leq 1 \text{ pour tt } i)
\]
est injectif (\struck{\ill{} qu'il est aussi surjectif} ce qui est vérifié dans le cas d'une $n$-alité) signifie ceci

(iii) $\forall s\in S$, $x\in R$, existe $s'\in S$ tel que $\{s,s'\}$ ne soit pas une arête (p.ex. $s'=s$) et $(s',x)\in R$ (N.B. en vertu de (ii) cet $s'$ est unique).
\marginal{\uncertain{implique que tous les $R(x)$ ont même cardinal $n = \mathrm{card}\,I$}}

Alors on a
\[
\varphi_i : E \longrightarrow E_i
\]
en associant : \struck{tout $x\in E$ l'unique} \struck{à} $\varphi_i(x)\in E_i$ tel que $\{\varphi_i(x)\} = S(x)\cap E_i$.

(iv) \struck{\ill{}} $\exists x\in E$ tel que $R(x)$ de card $\geq 2$. ($\Longleftrightarrow \mathrm{card}\,I\geq 2$)

\emph{N.B.} Dans le cas contraire, tous les $R(x)$ \uncertain{sont de} cardinal 1 (on suppose $S\neq\emptyset$) et $R$ est le graphe d'une application $E\to S$.

\section{Bi-icosaèdres et cartes de congruences}
\note{titre pris sur la feuille de garde (page 51 du fonds), de sa main : « Bi-icosaèdres et cartes de congruences, avril 1986 »}

\page{52}
\subsection{Bi-icosaèdres \add{\uncertain{gauches} et} orientations}
\marginal{Avril 86}

Soit $(S,\{\mathcal I,\mathcal I'\})$ un bi-icosaèdre \add{gauche} orienté, et soit $\widetilde S_{\mathcal I}, \widetilde A_{\mathcal I}, \widetilde F_{\mathcal I}$ le revêtement \struck{des} orienté de $(S,\mathcal I)$ \add{(i.e. ens. des icos. « droits » orientés)}, $\widetilde S_{\mathcal I'}, \widetilde A_{\mathcal I'}, \widetilde F_{\mathcal I'}$ celui de $(S,\mathcal I')$. (On va définir des bijections
\[
\begin{aligned}
\varepsilon^S_{\mathcal I',\mathcal I} &: \widetilde S_{\mathcal I} \xrightarrow{\ \sim\ } \widetilde S_{\mathcal I'}\\
\varepsilon^A_{\mathcal I',\mathcal I} &: \widetilde A_{\mathcal I} \xrightarrow{\ \sim\ } \widetilde A_{\mathcal I'}\\
\varepsilon^F_{\mathcal I',\mathcal I} &: \widetilde F_{\mathcal I} \xrightarrow{\ \sim\ } \widetilde F_{\mathcal I'}
\end{aligned}
\]
(pour \uncertain{vérification}, il faut \struck{\ill{}} remplacer $\widetilde A_{\mathcal I}, \widetilde A_{\mathcal I'}$ par $\overline A_{\mathcal I}$ et $\overline A_{\mathcal I'}$, ens. des arêtes \emph{orientées}) comme suit.

\struck{\ill{}} La fibre de $\widetilde S_{\mathcal I}$ et $\widetilde S_{\mathcal I'}$ en un point $s\in S$ sont les \add{\uncertain{deux} orientations des} deux \struck{structures} pentagonales sur $S(s)$ associées \uncertain{l'une} à $\mathcal I$, l'autre à $\mathcal I'$, et qui se \uncertain{correspondent}. \struck{\ill{}} \uncertain{Associé} au groupe des rotations $U_s$, \struck{isomorphe} qui est un \struck{espace aff.} $\mathbb F_5$-module libre de rang 1, et $U_s^*$ (ens. des générateurs) est un torseur sous \struck{\ill{}} le groupe cyclique $\mathbb F_5^*$ isomorphe à $\mathbb Z/4\mathbb Z$, \struck{\ill{}} ; $\widetilde S_{\mathcal I,s}$ est l'une des classes modulo le sous-groupe $\pm1$, soit $\{u,-u\}$, et $\widetilde S_{\mathcal I',s}$ est l'autre classe, donc $\{2u,-2u\}$.
\drawing{un pentagone avec son étoile de diagonales intérieures, en marge}
\marginal{On \uncertain{verrait} \ill{} \ill{} \ill{}}

L'application \add{\uncertain{i.e. dérivée}}
\[
\varepsilon^S_{\mathcal I',\mathcal I}(s) : \widetilde S_{\mathcal I}(s) \longrightarrow \widetilde S_{\mathcal I'}(s)
\]
est donnée par
\[
u \longmapsto 2u .
\]
On aura
\[
\varepsilon^S_{\mathcal I,\mathcal I'}(s)\,\varepsilon^S_{\mathcal I',\mathcal I}(s) = u\longmapsto 4u = -u = a_{\mathcal I'}u
\]
($a_{\mathcal I'}$ l'antipodisme de $\mathcal I'$), et de même en sens inverse. (Si on avait choisi $u\mapsto 3u$ ($=-2u$), ça aurait été pareil, car $3^2 = (-2)^2 \equiv -1\ (5)$) \struck{$4\equiv -1\ (5)$}. Donc
\[
\text{(1)}\qquad
\begin{cases}
\varepsilon_{\mathcal I',\mathcal I}\,\varepsilon_{\mathcal I,\mathcal I'} = a_{\mathcal I'}\\
\varepsilon_{\mathcal I,\mathcal I'}\,\varepsilon_{\mathcal I',\mathcal I} = a_{\mathcal I}
\end{cases}
\qquad\text{et de même.}
\]
\marginal{donc l'identification \uncertain{entre} $\widetilde S_{\mathcal I}$ et $\widetilde S_{\mathcal I'}$ dépend du choix de $\mathcal I$ ou de $\mathcal I'$ dans $\{\mathcal I,\mathcal I'\}$}
\marginal{On a \ill{} choisi $\mathcal I$ dans $\{\mathcal I,\mathcal I'\}$ \ill{} \ill{} $a_{\mathcal I}$ (ou $a_{\mathcal I'}$) \struck{\ill{}}}

Soit une face $t$ de $\mathcal I$, $t'$ sa complémentaire, face de $\mathcal I'$. On a une bijection canonique
\[
\alpha_{\mathcal I',\mathcal I} : t_{\mathcal I} \simeq t'_{\mathcal I'}
\]
(et $\alpha_{\mathcal I,\mathcal I'}$ est l'inverse !) d'où
\[
\varepsilon^F_{\mathcal I',\mathcal I} : \underset{\textstyle\widetilde A_{\mathcal I,t}}{\omega(t_{\mathcal I})} \xrightarrow{\ \sim\ } \underset{\textstyle\widetilde A_{\mathcal I',t'}}{\omega(t'_{\mathcal I'})}
\]
\note{sous chacun des deux membres, un signe d'égalité vertical renvoie à $\widetilde A_{\mathcal I,t}$ et $\widetilde A_{\mathcal I',t'}$ ; on attendrait $\widetilde F$, mais la lettre lue est $\widetilde A$}

\page{53}
d'où $\varepsilon^F_{\mathcal I',\mathcal I}$. Cette fois on aura
\[
\begin{cases}
\varepsilon^F_{\mathcal I',\mathcal I}\circ\varepsilon^F_{\mathcal I,\mathcal I'} = \mathrm{id}_{F_{\mathcal I'}}\\
\varepsilon^F_{\mathcal I,\mathcal I'}\circ\varepsilon^F_{\mathcal I',\mathcal I} = \mathrm{id}_{F_{\mathcal I}}
\end{cases}
\qquad\text{et de même.}
\]

Soit enfin $a$ une arête commune à $\mathcal I$, $\mathcal I'$, i.e. $a\in\mathfrak P_2(S)$. Alors $S\smallsetminus\{a\} = \alpha_{\mathcal I}\sqcup\alpha_{\mathcal I'}$, où $\alpha_{\mathcal I}$ est formé des deux éléments \struck{\ill{}} $u$ de $S\smallsetminus\{a\}$ tels que $(S\smallsetminus\{a\})\cup\{u\}\in\mathcal I$, et itou pour $\alpha_{\mathcal I'}$. On a bij. can.
\note{la lecture « $S\smallsetminus\{a\}$ » est celle de la page ; on attendrait, pour la seconde occurrence, la face $a\cup\{u\}$}
\drawing{un losange formé de deux triangles ayant l'arête $a$ en commun, avec une flèche courbe au-dessus : « carré $Q_{\mathcal I}$ »}
\[
\begin{aligned}
\widetilde A_{\mathcal I,a} &\simeq a\wedge\alpha_{\mathcal I} \quad (\simeq\omega(Q_{\mathcal I})) \quad\text{et de même}\\
\widetilde A_{\mathcal I',a} &\simeq a\wedge\alpha_{\mathcal I'}
\end{aligned}
\]
donc
\[
\widetilde A_{\mathcal I,a}\wedge\widetilde A_{\mathcal I',a} \simeq \alpha_{\mathcal I}\wedge\alpha_{\mathcal I'} \simeq a
\]
(N.B. on a $a\wedge\alpha_{\mathcal I}\wedge\alpha_{\mathcal I'} \overset{\text{can}}{=} 1$ à cause de l'orientation canonique des \ill{} d'un $S$ \uncertain{contenu} dans $\{\mathcal I,\mathcal I'\}$)

Donc choisir une bijection entre $\widetilde A_{\mathcal I,a}$ et $\widetilde A_{\mathcal I',a}$, \struck{\ill{}} i.e. choisir un \uncertain{élément} de $a$ \add{i.e. une orientation de l'arête $a$}. Donc on \ill{} \ill{} pas naturellement ici de bijection canonique au niveau de $\widetilde A_{\mathcal I}$ et $\widetilde A_{\mathcal I'}$ comme ci-dessus. Ainsi, on a une bijection
\[
\varepsilon^{\vec A}_{\mathcal I',\mathcal I} : \vec{\widetilde A}_{\mathcal I} \xrightarrow{\ \sim\ } \vec{\widetilde A}_{\mathcal I'}
\]
et de même en sens inverse, satisfaisant cette fois
\[
\begin{cases}
\varepsilon^{\vec A}_{\mathcal I',\mathcal I}\,\varepsilon^{\vec A}_{\mathcal I,\mathcal I'} = \mathrm{id}_{\vec{\widetilde A}_{\mathcal I'}}\\
\varepsilon^{\vec A}_{\mathcal I,\mathcal I'}\,\varepsilon^{\vec A}_{\mathcal I',\mathcal I} = \mathrm{id}_{\vec{\widetilde A}_{\mathcal I}}
\end{cases}
\qquad\text{et de même.}
\]
\marginal{N.B. $\vec{\widetilde A}_{\mathcal I}\simeq\mathrm{Rep}^+(S,A,\mathcal I)$, $\vec{\widetilde A}_{\mathcal I'}\simeq\mathrm{Rep}^+(S,A,\mathcal I')$}

En résumé : Pour les dodécaèdres \add{$D$, $D'$} \struck{des} \uncertain{duaux} des icosaèdres $(S,A,\mathcal I) = I$ et $(S,A,\mathcal I') = I'$, \struck{\ill{}} \uncertain{on peut} considérer $D$, $D'$ \uncertain{précisés} par la \ill{} \ill{} des sommets (le choix de l'identification ne dépendant pas d'un choix d'un $\mathcal I\in\{\mathcal I,\mathcal I'\}$), mais l'analogue \uncertain{pour} \uncertain{les icosaèdres} $I$, $I'$ n'est pas vrai : le choix d'une identification des sommets $S(I)\simeq S(I')$ dépend du choix de $I$ ou $I'$

\page{54}
et est remplacé par son produit par l'antipodisme (ce dernier, on le voit, kif kif) si on remplace $\mathcal I$ par $\mathcal I'$.

On appellera \emph{bi-icosaèdre orienté} \struck{un système formé de deux objets $\{I,I'\}$ de la \ill{}} \struck{de} un polyèdre comb. orienté \ill{} deux comp. connexes $I$, $I'$, qui sont des icosaèdres, et la donnée d'une \add{paire $\{\Theta,\Theta'\}$} \struck{classe de} \struck{deux} bijections \add{($=$ correspondances)} \struck{entre} $S(I)$ et $S(I')$ compatibles aux antipodismes, et \struck{\ill{}} \add{tel que} $\Theta'$ et $\Theta$ se déduisent l'un de l'autre par comp. avec l'antipodisme, et tel que, en prenant en quotient \add{de} $\Theta$ et $\Theta'$ modulo antipodismes, la bijection qu'on obtient entre $S(I)/\pm$ et $S(I')/\pm$ \struck{transforme la bijection \ill{}} \ill{} sur l'un comme une structure \uncertain{icosaédrale} \uncertain{gauche} \ill{} complémentaire !
\marginal{(et les bijections \ill{} \uncertain{définies} \ill{} par $\varepsilon^S_{\mathcal I,\mathcal I'}$ et $\varepsilon^S_{\mathcal I',\mathcal I}$)}

Il revient au même de dire que l'on se donne \struck{une donnée d'\ill{}} \add{une paire d'iso orientés}
\[
\begin{cases}
\vec I' \xrightarrow[\sim]{\ \alpha_{I'}\ } \mathrm{compl.}(\vec I)\\
\vec I \xrightarrow[\sim]{\ \alpha_I\ } \mathrm{compl.}(\vec I')
\end{cases}
\]
\struck{$I'$ avec l'icos. complémentaire} \marginal{icos. « complémentaire » de $\vec I$, avec l'orient. et les \uncertain{orientations} canoniques déduites}
dont les composés soient les antipodismes
\[
\mathrm{compl}(\alpha_{I'})\,\alpha_I : I \longrightarrow \underbrace{\mathrm{compl}\,\mathrm{compl}(I)}_{\textstyle -I}
\]
($-I$, i.e. $I$ avec l'orientation opposée) soit l'antipodisme, qui est un iso de $I$ avec $-I$.

\page{55}
Si on a deux bi-icosaèdres orientés $\{I_1,I'_1\}$ et $\{I_2,I'_2\}$, le choix d'une \struck{icos. orient.} composante connexe des premiers, et d'un iso de celle-ci avec une comp. conn. de l'autre, se prolonge de façon unique en un iso de bi-icosaèdres orientés.

La paire \struck{\ill{}} $(\widetilde I_{\mathcal I}, \widetilde I_{\mathcal I'})$ des rev. orientés des deux structures icosaédrales \add{gauches} \struck{complémentaires}, est un bi-icos. orienté, et son groupe des automorphismes est celui du \struck{groupe} bi-icosaèdre \uncertain{gauche} de départ.

Une autre façon de dire les choses \ill{} bi-icosaèdre droit \ill{} celle-ci. Considérons un icosaèdre orienté $I$. On lui associe l'icosaèdre complémentaire $I^\circ$, tel que pour $s,s'\in S = S(I)$, $\{s,s'\}$ soit arête de $I^\circ$ ssi $\{s,-s'\}$ est arête de $I$ (où $-$ est l'antipodisme). On a $I = I^{\circ\circ}$. \struck{On \ill{} le groupe des} Une structure de bi-icosaèdre sur $S$ (ensemble de douze sommets) est une paire de structures icosaédrales $\{I,I'\}$ sur $S$, mutuellement complémentaires. Le groupe des automorphismes s'envoie dans le groupe \add{(\uncertain{\ill{}}) le groupe des} permutations de $\{I,I'\}$, d'où
\[
1 \longrightarrow \underset{\substack{\simeq\,\mathrm{Aut}(I')\\ \simeq\,\mathrm{Aut}^+(I)\times\pm1\\ \simeq\,\mathrm{Aut}^+(I')\times\pm1}}{\mathrm{Aut}(I)} \longrightarrow \mathrm{Aut}(I,I') \longrightarrow \mathbb Z/2\mathbb Z \longrightarrow 1
\]

\page{56}
Le groupe quotient
\[
\mathrm{Aut}(I,I')\,/\,\underbrace{\mathrm{Aut}^{0+}(I,I')}_{\substack{\text{automorphismes qui respectent } I,\ I'\\ \text{et l'orientation de } I \text{, i.e. de } I'}}
\]
est isomorphe à $\underline{\mathbb Z/4\mathbb Z}$ sauf erreur. Mais si on veut se donner une \struck{\ill{} bi-icosaèdre orienté} orientation de $I$ et chacun de $I'$ (pour définir une notion de bi-icosaèdre \emph{orienté}) \uncertain{définie} une notion \ill{} $\vec I$ et $\vec I'$, \struck{\ill{}} on ne peut pas \struck{\ill{}} \add{d'\ill{}} de choisir \struck{\ill{}} une des icos. de la structure bi-icosaédrale de façon canonique — comme celui pour laquelle l'orientation de l'autre soit « induite ». Donc on ne trouve pas la bonne notion !

Par contre, il y a un groupe $\mathbb Z/4\mathbb Z$ qui opère sur la catégorie des icosaèdres orientés, dont le générateur est \struck{le passage \ill{}} le passage de $\vec I$ à $\vec I^\circ$. (Le carré du générateur est le passage de $I$ à $-I$ — correspondant à l'ext. triviale de $\mathbb Z/2\mathbb Z\subset\mathbb Z/4\mathbb Z$ par $\mathrm{Aut}(\vec I)^+$, savoir le groupe des automorphismes de \struck{$\mathrm{Aut}$} $I$ \uncertain{lui} respectant l'orientation)

\page{57}
Ainsi, il y a deux extensions, \uncertain{à iso près} \uncertain{indépendantes}, \ill{} comme \ill{} du groupe de l'icosaèdre orienté
\[
\begin{gathered}
0 \longrightarrow \mathrm{Aut}(\vec I) \longrightarrow \mathrm{Aut}(\underbrace{\{\vec I,\vec I^\circ\}}_{\text{bicos. or.}}) \longrightarrow \mathbb Z/2\mathbb Z \longrightarrow 0\\
0 \longrightarrow \mathrm{Aut}(\vec I) \longrightarrow \mathrm{Aut}(\{I,I^\circ\}) \longrightarrow \mathbb Z/4\mathbb Z \longrightarrow 0
\end{gathered}
\]

Je n'y \uncertain{pige} pas grand-chose, d'ailleurs dans un sens \uncertain{ou} dans l'autre, i.e. \uncertain{lorsque} l'identité de $\mathrm{Aut}(\vec I)$. De haut en bas, le gén. de $\mathbb Z/2\mathbb Z$ (qui définit la classe d'ext. triviale de $\mathrm{Aut}(\vec I)$) devrait aller sur l'un des deux gén. de $\mathbb Z/4\mathbb Z$, impossible ! Dans l'autre sens, un \struck{\ill{}} $\mathbb Z/4\mathbb Z\to\mathbb Z/2\mathbb Z$ devrait être surjectif (\uncertain{ça va}) et le noyau \struck{devrait} dans $\mathrm{Aut}(I,I')$ est un sous-groupe d'ordre 2 — l'antipodisme ! Et en fait, je vois maintenant que $\mathrm{Aut}(\vec I,\vec I^\circ) \xrightarrow{\sim}$ \struck{le quotient} au quotient de $\mathrm{Aut}(I,I')$ par l'antipodisme. Mais cela semblerait indiquer qu'il y aurait un foncteur
\[
\text{bi-icosaèdres (droits) non orienté} \longmapsto \text{bi-icosaèdre orienté}
\]
et en effet, on passe par bi-icosaèdres gauches, et bi-icosaèdres orientés via rev. des orientations.

\page{58}
La théorie générale des cartes de congruence nous donne
\[
1 \longrightarrow \underbrace{\mathrm{Sl}(2,\mathbb F_5)/\pm1}_{\mathrm{Aut}(\vec I)} \longrightarrow \mathrm{Gl}(2,\mathbb F_5)/\pm1 \longrightarrow \underset{\textstyle\mathbb Z/4\mathbb Z\ !!!}{\mathbb F_5^*} \longrightarrow 1
\]
\note{sous $\mathbb F_5^*$, un signe $\simeq$ vertical}
Le groupe $\mathrm{Gl}(2,\mathbb F_5)$ permute entre elles l'une des 4 structures d'icos. orienté \uncertain{portées} par l'ens. $S$ à 12 éléments, et on \uncertain{voit} le groupe cyclique $\mathbb Z/4\mathbb Z$.

On doit pouvoir plonger les deux extensions (de $\mathbb Z/4\mathbb Z$, et de $\mathbb Z/2\mathbb Z$) par $\mathrm{Aut}^+(\vec I^\circ)$, dans le groupe produit \struck{\ill{}} \uncertain{$\mathrm{PGl}(2,\mathbb F_5)$} $\times\,\mathbb F_5^*$ (\uncertain{$\simeq\mathfrak S_5$, ordre 120}), \uncertain{biunivoquement} (et \uncertain{canoniquement}) dans
\[
G = \mathrm{Aut}(I)\times\underset{\textstyle\simeq\mathbb Z/4\mathbb Z}{\mathbb F_5^*} .
\]
Ce groupe est réalisé comme un groupe d'automorphismes des \struck{polyèdre} \add{doubles} \struck{\ill{} bi-icosaèdre} $I\sqcup I'$, mais en respectant la structure supplémentaire donnée par un iso $I/\text{antip.}\simeq I'/\text{antip.}$ ($\simeq J$ \uncertain{i.c.g.}) ! Choisissons une orientation sur $I$, une autre sur $I'$ (ce \add{qui} \uncertain{fournit} d'ailleurs $\vec I$, $\vec I'$ \uncertain{relevant} \add{\uncertain{deux} \ill{}} au rev. \struck{universel} \add{orienté} de $J$) on trouve le sous-groupe respectant ces orientations : $\mathrm{Aut}^+(I) = \mathrm{Aut}^+(\vec I)$. Le quotient \add{(d'ordre 8)} opère sur l'ens. des quatre orientations de $(I,I')$, et \struck{les}
\marginal{N.B. ce nouveau groupe \ill{} \ill{} \ill{} groupe \ill{} droits \ill{} orientations}

\page{59}
c'est le groupe des automorphismes du carré combin. \uncertain{orienté} \struck{permute entre elles \ill{} ($\mathfrak S_5$), \ill{} \ill{} \ill{} \ill{}}. L'ext. de $\mathbb Z/2\mathbb Z$ par $\mathrm{Aut}^+(\vec I)$ est l'image inverse du \struck{\ill{}} sous-groupe d'ordre 2 \struck{des $\{\pm1\}$ de $\mathbb F_5^*$ (\ill{} de l'antipodisme des côtés)} correspondant au choix des deux orientations (mais qui \struck{permute les \ill{} des systèmes, respectant la paire d'orientations (\ill{} \ill{})} \add{permute} ses deux extrémités). L'ext. de $\mathbb Z/4\mathbb Z$ par $\mathrm{Aut}^+(\vec I)$ s'obtient en prenant l'image inverse du groupe des rotations du carré. L'intersection de ces deux sous-groupes de $G$ est $\mathrm{Aut}^+(\vec I)$ (\uncertain{\ill{}} le groupe des automorphismes du carré). Le groupe engendré est formé des produits, et c'est $G$ tout entier…

\emph{Généralisation} aux cartes de congruence

La cat. $\mathrm{Congr}$ est réalisée avec \add{\uncertain{l'ens. des sommets}} \add{\uncertain{primitifs}} \struck{sommets qui sont des} paires $\{a,-a\}$ d'éléments d'un $\mathbb Z/n\mathbb Z$-module \add{symplectique $E$} libre de rang 2, formant base, on prend
\[
a\otimes a \in \mathrm{End}(E) = M_2(\mathbb Z/n\mathbb Z)
\]
\uncertain{qui sont} des éléments du cône isotrope $C$ de \uncertain{$M^0$} (éléments de trace nulle) i.e. formé des nilpotents. On trouve ainsi une réalisation
\marginal{\ill{} \ill{} si $\mathbb Z/n\mathbb Z$ \ill{} corps, \ill{} ?}

\page{60}
(compatible avec les op. du groupe $G_n = $ \struck{$\mathrm{Sl}(2,\mathbb Z)$} $\mathrm{Sl}(E)/\pm1$, mais pas du groupe $\mathrm{Gl}(E)$ tout entier) de $S$ dans ledit cône. On ne trouve pas tous les éléments non dégénérés dudit cône, mais \struck{seulement \ill{}} \add{ceux \uncertain{dans} une \uncertain{orbite}} \add{$\mathcal I_\lambda$ ($\lambda\in\mathbb F_p^*$)} classe mod $\mathbb F_p^*/\mathbb F_p^{*2}$. Les opérations de $\mathbb F_p^* = \mathbb Z/n\mathbb Z^*$ sur $S$, se transforment en les opérations d'homothétie par les $\lambda^2$. Mais les homothéties par les $\lambda\in\mathbb F_p^*\smallsetminus\mathbb F_p^{*2}$ échangent entre elles les deux composantes connexes de l'un des \struck{pts} \add{nuls} \struck{\ill{}} des cônes (cette division en classes, \struck{indexées par} \add{\uncertain{liée}} à $k$, l'une des classes formant un torseur sous $k^*/k^{*2}$, a un sens pour tt alg. d'Azumaya de rang 4 sur un corps quelconque $k$ — si $k$ est fini \struck{\ill{}} $\mathbb F_q$ de car $\neq 2$, $k^*/k^{*2}$ est cyclique d'ordre pair, $k^*/k^{*2}$ est d'ordre 2 …) Ainsi on trouve une « bi-structure » \uncertain{par} quotient de \uncertain{$\mathrm{Gl}(E)$} $\times\,\mathbb F_p^*$ dans un…

TSVP
\note{l'argument continue au-delà de ce lot}

\end{document}
